\section{训练监控与日志体系}

\subsection{日志的结构化写法}

复制整个训练过程的关键是详细的 logging schema。以下字段是最低要求：

\begin{itemize}
    \item `step`, `episode`, `wall_time`：定位发生问题的 slice
    \item `avg_return`, `std_return`, `kl_divergence`, `entropy`：衡量 policy 未来趋势
    \item `buffer_size`, `td_error_mean`, `td_error_std`：Replay buffer 内部状态
    \item `learning_rate`, `grad_norm`, `value_loss`：诊断优化器与 value head
\end{itemize}

\begin{importantbox}{日志粒度与报警}
高频 logging（每 1k step）结合 threshold-based alarms（KL > 0.06, return variance > 0.2）能在 divergence 发生前打断训练，特别是在 Sim2Real 环境中避免一次物理损坏。
\end{importantbox}

\subsection{指标面板}

实时可视化 dashboard 通常至少包含：

\begin{enumerate}
    \item PPO ratio distribution histogram
    \item SAC entropy & Q-value curves
    \item Batch-wise advantage histogram to detect saturation
\end{enumerate}

\subsection{本章小结}

结构化 logging 与实时 dashboard 能让研究者及时发现 divergence、explosion、overfitting 等 failure modes，从而更快复现核心表现。

